% theme: quadratic_equations_percent_applications
\MajorQuestionLesson{割合の利用問題}
\SetRowSpace{6mm}

\TypeQuestionSingle{次の問いに答えなさい。}{kingaku}
\QQ{2500円で仕入れた商品に $x$\percent の利益を見込んで定価をつけた。売れ残ったため、定価の $x$\percent 引きで売ったところ、仕入れ値より225円だけ安くなった。$x$ の値を求めなさい。}{$x=30$}
\QQ{定価4000円の商品を $x$\percent 引きにし、さらにその値段から $x$\percent 引きにしたところ、値段は2560円になった。$x$ の値を求めなさい。}{$x=20$}
\QQ{定価8000円の商品を $x$\percent 引きにし、さらにその値段から $(x+10)$\percent 引きにしたところ、値段は4480円になった。$x$ の値を求めなさい。}{$x=20$}
\QQ{2000円で仕入れた商品に $x$割の利益を見込んで定価をつけた。売れ残ったため、定価の $3x$割引きで売ったところ、売値は1540円になった。$x$ の値を求めなさい。}{$x=1$}
\QQ{定価9000円の商品を $x$割引きにし、さらにその値段から $(x+1)$割引きにしたところ、値段は5040円になった。$x$ の値を求めなさい。}{$x=2$}

\TypeQuestionSingle{次の問いに答えなさい。}{ninzuu}
\QQ{ある学校の生徒数は1000人であった。生徒数が毎年同じ割合で増え、2年後に1210人になった。1年あたり何\percent 増えたか求めなさい。}{$10$\percent}
\QQ{ある町の人口は1000人であった。人口が毎年同じ割合で減り、2年後に810人になった。1年あたり何\percent 減ったか求めなさい。}{$10$\percent}
\QQ{ある学校の生徒数は1000人であった。1年目に $x$\percent 増え、2年目に $(x+10)$\percent 増えたところ、生徒数は1320人になった。$x$ の値を求めなさい。}{$x=10$}
\QQ{ある町の人口は1200人であった。1年目に $x$\percent 減り、2年目に $(x+5)$\percent 減ったところ、人口は918人になった。$x$ の値を求めなさい。}{$x=10$}
\QQ{ある学校の男子は900人、女子は400人であった。男子は毎年 $x$\percent 減り、女子は毎年 $x$\percent 増えたところ、2年後に男子と女子の人数が同じになった。$x$ の値を求めなさい。}{$x=20$}

\LessonPageBreak

\TypeQuestionSingle{次の問いに答えなさい。}{noudo}
\QQ{$x$ gの $x$\percent の食塩水に60gの水を加えたところ、濃度が16\percent になった。$x$ の値を求めなさい。}{$x=40$}
\QQ{$(x+20)$gの $x$\percent の食塩水に20gの水を加えたところ、濃度が30\percent になった。$x$ の値を求めなさい。}{$x=40$}
\QQ{160gの10\percent の食塩水に $x$ gの食塩を加えたところ、濃度が $x$\percent になった。$x$ の値を求めなさい。}{$x=20$}
\QQ{$x$ gの $x$\percent の食塩水と80gの10\percent の食塩水を混ぜたところ、濃度が20\percent になった。$x$ の値を求めなさい。}{$x=40$}
\QQ{$x$ gの $x$\percent の食塩水に50gの水を加えたところ、濃度が25\percent になった。$x$ の値を求めなさい。}{$x=50$}

\TypeQuestionSingle{次の問いに答えなさい。}{shigoto}
\QQ{ポンプの1分あたりに入れる水の量を $x$\percent 増やし、水を入れる時間を $x$\percent 短くしたところ、入った水の量はもとの96\percent になった。$x$ の値を求めなさい。}{$x=20$}
\QQ{ポンプの1分あたりに入れる水の量を $x$\percent 増やし、水を入れる時間を $(x+10)$\percent 短くしたところ、入った水の量はもとの88\percent になった。$x$ の値を求めなさい。}{$x=10$}
\QQ{ある機械の1時間あたりの作業量を $x$\percent 増やし、作業時間を $(x+20)$\percent 短くしたところ、全体の作業量はもとの77\percent になった。$x$ の値を求めなさい。}{$x=10$}
\QQ{印刷機の1分あたりの印刷枚数を $x$\percent 増やし、印刷する時間を $x$\percent 短くしたところ、印刷した枚数はもとの99\percent になった。$x$ の値を求めなさい。}{$x=10$}
\QQ{ポンプの1分あたりに入れる水の量を $x$\percent 減らし、水を入れる時間を $(x+20)$\percent 長くしたところ、入った水の量はもとの117\percent になった。$x$ の値を求めなさい。}{$x=10$}
